\documentclass[10pt, a4paper]{article}

% Parametri che modificano il file main.tex
% Le uniche parti da cambiare su main.tex sono:
% - tabella versioni
% - testo file

\def\titolo{Analisi dei \\ \ requisiti} % \\ per andare a capo, \ per spazio
\def\titoloHeader{Analisi dei requisiti}
\def\data{}

% \def\listaComponenti{
% Bresolin G.,
% Campese M.,
% Ciriolo I.,
% Dugo A.,
% Feltrin E.,
% Michelon R.,
% Orlandi G.
% }

\usepackage{style}
\usepackage{headerfooter}

\title{\titolo}
\author{SWEetCode}

\begin{document}

% PRIMA PAGINA
\begin{titlepage}
    \thispagestyle{empty}
    \begin{tikzpicture}[remember picture, overlay]
        % TRIANGOLI
        \draw[fill=secondarycolor, secondarycolor] (current page.north west) -- (current page.south west) -- (8.8, -28);
        \draw[fill=primarycolor, primarycolor] (-3, 5) -- (4, -13.6) -- (11, 5);

        % LOGO
        \node [xshift=-5cm, yshift=25cm] (logo) at (current page.south east) {\includegraphics[width=6.5cm]{img/logo.png}};

        % SWEETCODE - DATE
        \node [anchor=north east, align=right, xshift=-1.2cm, yshift=20.5cm, text=black] (sweetcode) at (current page.south east) {\fontsize{32pt}{36pt}\selectfont SWEetCode};
        \draw[line width=4pt, lightcol] ([xshift=-3cm, yshift=-0.37cm]sweetcode.south west) -- ([yshift=-0.37cm]sweetcode.south east);
        \node [anchor=north east, align=right, xshift=-1.2cm, yshift=18.7cm, text=black] (date) at (current page.south east){\fontsize{24pt}{24pt} \selectfont \data};

        % NOME FILE
        \node [anchor=north east, text width=15cm, align=right, xshift=-1.2cm, yshift=17cm, text=black] (titolo) at (current page.south east){\fontsize{48pt}{48pt}\textbf{\titolo}};

        % BOX DATI PARTECIPANTI
        \node[anchor=north east, xshift=-1.2cm, yshift=12cm, minimum width=8cm] (box) at (current page.south east){};

        % COMPONENTI DEL GRUPPO
        \node[anchor=north west, yshift=-1cm, align=left] (dati1) at (box.north west) {\fontsize{16pt}{16pt}\selectfont \textbf{Componenti del gruppo}};
        \draw[line width=4pt, lightcol] (dati1.south west) -- ([xshift=8cm]dati1.south west);
        
        \node[anchor=north west, align=left] (dati11) at (dati1.south west)
        {\fontsize{14pt}{14pt}\selectfont Bresolin G.};
        \node[anchor=north west, yshift=-0.7cm, align=left] (dati12) at (dati1.south west)
        {\fontsize{14pt}{14pt}\selectfont Campese M.};
        \node[anchor=north west, yshift=-1.4cm, align=left] (dati13) at (dati1.south west)
        {\fontsize{14pt}{14pt}\selectfont Ciriolo I.};
        \node[anchor=north west, yshift=-2.1cm, align=left] (dati14) at (dati1.south west)
        {\fontsize{14pt}{14pt}\selectfont Dugo A.};
        \node[anchor=north west, yshift=-2.8cm, align=left] (dati15) at (dati1.south west)
        {\fontsize{14pt}{14pt}\selectfont Feltrin E.};
        \node[anchor=north west, yshift=-3.5cm, align=left] (dati16) at (dati1.south west)
        {\fontsize{14pt}{14pt}\selectfont Michelon R.};
        \node[anchor=north west, yshift=-4.2cm, align=left] (dati17) at (dati1.south west)
        {\fontsize{14pt}{14pt}\selectfont Orlandi G.};
        
        % UNIPD - SWE
        \node [xshift=4.4cm, yshift=2.3cm, draw, secondarycolor, text=white] (uni) at (current page.south west) {\fontsize{20pt}{20pt} \selectfont Università di Padova};
        \node [xshift=0.65cm, yshift=0.7cm, draw, secondarycolor, text=white, below=of uni] (corso) {\fontsize{20pt}{20pt}\selectfont Ingegneria del Software};

        % FIRMA
        % \draw[line width=4pt, lightcol] ([xshift=-1.2cm, yshift=1.8cm]current page.south east) -- ([xshift=-8cm, yshift=1.8cm]current page.south east);
        % \node[anchor=north west, xshift=12.9cm, yshift=1.45cm, align=left] at (current page.south west)
        % {\fontsize{13pt}{13pt}\selectfont L'Amministratore: Feltrin E.};
        
    \end{tikzpicture}
\end{titlepage}

% REGISTRO DELLE VERSIONI
%Registro in ordine dalla più recente alla meno recente!
{\renewcommand{\arraystretch}{1.5}
\section*{Registro delle versioni}

\begin{xltabular}{\textwidth}{c|c|c|c|X}
\label{tab:long}

\textbf{Versione} & \textbf{Data} & \quantities{\textbf{Responsabile di}\\\textbf{stesura}}& \textbf{Revisore} & \quantities{\textbf{Dettaglio e}\\\textbf{motivazioni}} \\
\endfirsthead

\textbf{Versione} & \textbf{Data} & \quantities{\textbf{Responsabile di}\\\textbf{stesura}}& \textbf{Revisore} & \quantities{\textbf{Dettaglio e}\\\textbf{motivazioni}} \\
\endhead

\multicolumn{5}{r}{{Continua nella pagina successiva}} \\
\endfoot

\endlastfoot

\hline
v2.0.0(1) & $2023-11-20$ & \quantities{Bresolin G.\\Campese M.\\Dugo A.\\Michelon R.\\Feltrin E.\\Ciriolo I.} & Orlandi G. & Aggiornata sezione glossario.\\
\hline
v1.0.0(1) & $2023-11-19$ & Michelon R. & Campese M. & Modificata introduzione.\\
\hline
v1.0.1 & $2023-11-14$ & \quantities{Ciriolo I.} & Campese M. &  Studio dei primi casi d'uso.\\
\hline
v0.0.1 & $2023-11-06$ & \quantities{Ciriolo I.\\Orlandi G.} & Campese M. &  Impostazione del documento, Introduzione.\\
    
\end{xltabular}}
\newpage

% INDICE
\tableofcontents
\newpage

% INTRODUZIONE
\section{Introduzione}
\subsection{Scopo del documento}
Il documento ha lo scopo di definire i casi d'uso dell'applicazione che verrà sviluppata nel corso del progetto. Vengono inoltre presentati i requisiti funzionali, di qualità e di vincolo e le funzionalità aggiuntive che saranno implementate non essendo ritenute essenziali/necessarie.
\\
%\subsection{Scopo del prodotto}
\subsection{Glossario}
Per evitare ambiguità e incomprensioni relative al linguaggio e ai termini utilizzati nella documentazione relativa al progetto viene presentato un Glossario. I termini ambigui o specifici presenti nello stesso, verranno identificati con un pedice |g|. (DA VALUTARE COME RICONOSCERE IL TERMINE )
\\
\subsection{Riferimenti}
   \subsubsection{Riferimenti normativi}
   \begin{itemize}
    \item \textit{Regolamento del progetto didattico}: \\
    \href{https://www.math.unipd.it/~tullio/IS-1/2023/Dispense/PD2.pdf}{https://www.math.unipd.it/~tullio/IS-1/2023/Dispense/PD2.pdf};
    \item \textit{Norme di progetto v0.4.3}.
    \end{itemize}
    
    \subsubsection{Riferimenti informativi}
    
    \begin{itemize}
    \item \textit{Glossario v0.0.1} (da creare parallelamente); 
    \item \textit{Presentazione capitolato C1}:\\
    \href{https://www.math.unipd.it/~tullio/IS-1/2023/Progetto/C1.pdf}{https://www.math.unipd.it/~tullio/IS-1/2023/Progetto/C1.pdf};
    \item \textit{Verbali esterni ed interni};
    \item \textit{Analisi e descrizione delle funzionalità: Use Case e relativi diagrammi (UML):}\\
    \href{https://www.math.unipd.it/~rcardin/swea/2022/Diagrammi%20Use%20Case.pdf}{https://www.math.unipd.it/~rcardin/swea/2022/Diagrammi\%20Use\%20Case.pdf ;}
    \item \textit{Analisi e descrizione delle funzionalità: Diagrammi delle Attività (UML):}\\
    \href{https://www.math.unipd.it/~rcardin/swea/2022/Diagrammi%20di%20Attivit%C3%A0.pdf}{https://www.math.unipd.it/~rcardin/swea/2022/Diagrammi\%20di\%20Attivit\%C3\%A0.pdf;}
    

    \item \textit{Progettazione e Programmazione: Diagrammi delle Classi:}\\
    \href{https://www.math.unipd.it/~rcardin/swea/2023/Diagrammi%20delle%20Classi.pdf}{https://www.math.unipd.it/~rcardin/swea/2023/Diagrammi\%20delle\%20Classi.pdf.}
    \end{itemize}

% DESCRIZIONE
\newpage
\section{Descrizione}
\subsection{Obiettivi del prodotto}
Il proponente richiede la creazione di un sito web , dotato di un componente di intelligenza artificiale , che sia d'aiuto a tutto il personale dell'azienda .
In particolare , il sistema deve rendere accessibili le informazioni, le regolamentazioni e le indicazioni aziendali ai dipendenti.
Deve fornire nuove modalità di formazione e assistenza sul lavoro .
Deve poter superare i limiti di un classico repository documentale . Rendendo le informazioni più fruibili , concentrandosi più sui contenuti , che sulla struttura .
Il sistema deve riuscire a comunicare con i lavoratori in modo semplice , qualsiaisi persona , anche la meno esperta ,riguardo moderne tecnologie , deve riuscire a comunicare con esso .
Infine , deve riuscire a migliorare il comportamento dei lavoratori , facendo in modo che rispettino le regole .
\subsection{Funzioni del prodotto}
Il software deve essere in grado di ricevere in input documenti di qualsiasi natura , in particolare pdf , documenti fatti con word , con la quale possa rispondere a qualsiasi domanda che gli sottoponga il lavoratore . 
\subsection{Caratteristiche del prodotto}
Il sistema deve poter aver accesso a tutti i file che riguardino la produttività aziendale , non può rispondere simultaneamente a più utenti , può gestirne uno alla volta e deve poter conoscere l'italiano e l'inglese . I documenti caricati devono poter rimanere accessibili ed inoltre deve esserci una storia della comunicazione tra ogni lavoratore ed esso . Inoltre deve aver meccanismi riguardanti l'autenticazione di un generico lavoratore .
\subsection{Utenti}
Gli utenti possono interagire con il sistema unicammente attraverso un'interfaccia e non possono assolutamente andare a vedere i file che il sistema gestisce . L'utente può chiedere informazioni riguardo dove si trovino le informazioni che gli interessano . Se l'utente chiede al sistema informazioni , non pertinenti , il sistema può e deve farlo notare .
% CASI D'USO
\newpage
\section{Casi d'uso}

\subsection{Obiettivi}
La seguente sezione ha l'obiettivo di presentare i diversi Casi d'Uso relativi all'applicazione \textit{Knowledge Managment AI}. Essi sono indicati con la notazione seguente: \\ UC.[X].[Y] \ \  in cui:
\begin{itemize}
\item UC sta per \textit{Use Case};
\item .[X] indica il numero del Caso d'Uso, presentati secondo successione gerarchica;
\item .[Y] indica il possibile numero del sotto-Caso d'uso, che aggiunge informazioni al Caso d'Uso [X]. Anche i sotto-Casi d'Uso sono presentati secondo successione gerarchica.
\end{itemize}

\subsection{Attori}
L'applicazione \textit{Knowledge Managment AI} richiede un singolo attore che rappresenta l'utente (ovvero il dipendente aziendale) che interagisce con il ChatBot. \\ Si specifica che l'azienda proponente non richiede l'autenticazione degli utenti nell'applicazione, poichè rappresenterebbe un rallentamento alla divulgazione rapida e semplice di documenti all'interno della rete aziendale; Per questo motivo l'autenticazione non viene presentata tra i requisiti.
\subsection{Lista dei casi d'uso}
%\begin{itemize}
    %\item 
%\end{itemize}

\renewcommand{\arraystretch}{1.5}
\begin{tabularx}{\textwidth}{|c|X|c|c|c|}\hline
\textbf{Caso d'uso} & \textbf{Attori} & \textbf{Precondizioni} & \textbf{Postcondizioni} & \textbf{Scenario principale} \\
\end{tabularx}


% REQUISITI
\newpage
\section{Requisiti}

\subsection{Requisiti funzionali}
Di seguito la specifica per i requisiti funzionali, i quali descrivono le funzionalità del sistema, le azioni
che il sistema può compiere e le informazioni che il sistema può fornire.
Si intende che la numerologia di ciascuno rispecchi ove previsto gli UC presenti, mentre le sigle sotto
riportate indicano:
\begin{itemize}
    \item ROF - Requisito Obbligatorio Funzionale;
    \item ROF - Requisito Obbligatorio Funzionale.
\end{itemize}
\renewcommand{\arraystretch}{1.5}
\begin{tabularx}{\textwidth}{|c|X|c|c|}\hline
\textbf{Codice} & \textbf{Descrizione} & \textbf{Classificazione} & \textbf{Fonti} \\
\end{tabularx}

\subsection{Requisiti di qualità}
Di seguito la specifica per i requisiti di qualità, i quali descrivono come un sistema deve essere, o
come il sistema deve esibirsi, per soddisfare le esigenze dell'utente.
Le sigle sotto riportate possono essere così classificate:
\begin{itemize}
    \item RDQ - Requisito Desiderabile di Qualità;
    \item ROQ - Requisito Obbligatorio di Qualità.
\end{itemize}
\renewcommand{\arraystretch}{1.5}
\begin{tabularx}{\textwidth}{|c|X|c|c|}\hline
\textbf{Codice} & \textbf{Descrizione} & \textbf{Classificazione} & \textbf{Fonti} \\
\end{tabularx}

\subsection{Requisiti di vincolo}
Di seguito la specifica per i requisiti di vincolo, i quali descrivono i limiti e le restrizioni che un sistema
deve rispettare per soddisfare le esigenze dell'utente.
Le sigle sotto riportate possono essere così classificate:
\begin{itemize}
    \item RDV - Requisito Desiderabile di Vincolo;
    \item ROV - Requisito Obbligatorio di Vincolo.
\end{itemize}
\renewcommand{\arraystretch}{1.5}
\begin{tabularx}{\textwidth}{|c|X|c|c|}\hline
\textbf{Codice} & \textbf{Descrizione} & \textbf{Classificazione} & \textbf{Fonti} \\
\end{tabularx}




\end{document}
